\chapter{Type I singularities}

\section*{1. Intuition}

\begin{remark}[Singularity types in dimension three]
\label{ch09-rem-singularity-intuition}
\source{chow-knopf-ricci-flow:page-0266, chow-knopf-ricci-flow:page-0267}
\dcref{Section 9.1}
\group{intuition}

In dimension three, finite-time singularities are expected to be tractable because of pinching, the structure of the curvature operator, and the strong maximum principle for tensors. A Type I singularity should resemble either a shrinking round sphere or a shrinking round cylinder, while a Type IIa singularity should lead after suitable dilations to one of the noncompact limits described in the source.
\end{remark}

\begin{remark}[Canonical model geometries]
\label{ch09-rem-canonical-model-geometries}
\source{chow-knopf-ricci-flow:page-0267, chow-knopf-ricci-flow:page-0268}
\dcref{Equation 9.1, Equation 9.2}
\group{intuition}

The standard Type I model in dimension three is the shrinking round sphere with metric \(4(T-t)\,g_{\mathrm{can}}\). The standard cylindrical model is the shrinking round product cylinder \(\mathrm{d}s^2+2(T-t)\,g_{\mathrm{can}}\) on \(\mathbb{R}\times S^2\).
\end{remark}

\begin{remark}[Heuristic list of singularity models]
\label{ch09-rem-singularity-model-heuristics}
\source{chow-knopf-ricci-flow:page-0267, chow-knopf-ricci-flow:page-0268}
\dcref{Section 9.1}
\group{intuition}

The source organizes the expected limits after dilation into the compact spherical case, the cylindrical case, and the noncompact positively curved steady or cigar-type cases that can arise after further dimension reduction.
\end{remark}

\section*{2. Positive curvature is preserved}

\begin{lemma}[Positive sectional curvature is preserved]
\label{ch09-lem-positive-sectional-curvature-preserved}
\source{chow-knopf-ricci-flow:page-0268, chow-knopf-ricci-flow:page-0269}
\dcref{Lemma 9.2}
\group{positive-curvature-is-preserved}

If \((\mathcal M^3,g(t))\) is a solution of the Ricci flow such that \(g(0)\) has positive, respectively nonnegative, sectional curvature, then \(g(t)\) has positive, respectively nonnegative, sectional curvature for all \(t>0\) for which the solution exists.
\end{lemma}

\begin{proof}
\source{chow-knopf-ricci-flow:page-0268, chow-knopf-ricci-flow:page-0269}
\sketch The curvature operator is encoded by a quadratic form whose smallest eigenvalue is preserved in the cone of positive sectional curvature by the ODE for the curvature operator. The key point is that the cone is fiberwise convex and invariant under parallel translation, and the evolution of the smallest eigenvalue is strictly inward whenever it is positive. The maximum principle for systems then implies preservation of positivity, and the nonnegative case is analogous.
\end{proof}

\begin{lemma}[Positive Ricci curvature is preserved]
\label{ch09-lem-positive-ricci-curvature-preserved}
\source{chow-knopf-ricci-flow:page-0269}
\dcref{Lemma 9.3}
\group{positive-curvature-is-preserved}

If \((\mathcal M^3,g(t))\) is a solution of the Ricci flow such that \(g(0)\) has positive, respectively nonnegative, Ricci curvature, then \(g(t)\) has positive, respectively nonnegative, Ricci curvature for all \(t>0\) for which the solution exists.
\end{lemma}

\begin{proof}
\source{chow-knopf-ricci-flow:page-0269}
\sketch The proof is the same maximum-principle argument as for sectional curvature, applied to the cone of curvature operators whose two smallest eigenvalues sum to a positive, respectively nonnegative, quantity.
\end{proof}

\section*{3. Positive sectional curvature dominates}

\begin{theorem}[Pinching estimate for the smallest curvature eigenvalue]
\label{ch09-thm-pinching-estimate}
\source{chow-knopf-ricci-flow:page-0271, chow-knopf-ricci-flow:page-0272, chow-knopf-ricci-flow:page-0273}
\dcref{Theorem 9.4}
\group{positive-section-dominates}

Let \((\mathcal M^3,g(t))\) be any solution of the Ricci flow on a closed \(3\)-manifold for \(0\le t<T\). Let \(\nu(x,t)\) denote the smallest eigenvalue of the curvature operator. If \(\inf_{x\in\mathcal M^3}\nu(x,0)\ge -1\), then at any point \((x,t)\in\mathcal M^3\times[0,T)\) where \(\nu(x,t)<0\), the scalar curvature satisfies
\[
R\ge |\nu|\bigl(\log |\nu|+\log(1+t)-3\bigr).
\]
\end{theorem}

\begin{proof}
\source{chow-knopf-ricci-flow:page-0271, chow-knopf-ricci-flow:page-0272, chow-knopf-ricci-flow:page-0273}
\sketch The source first establishes two convexity lemmas: a time-dependent convex cone of curvature operators and a differential inequality for the quantity \(\Omega=\tr(\mathbb P)/(-\nu(\mathbb P))-\log(-\nu(\mathbb P))\). These are then combined with the maximum principle for systems to show that the curvature operator remains in the cone for as long as \(\nu<0\). The claimed scalar-curvature lower bound follows by rewriting the cone condition in terms of \(R=\tr \Rm\) and \(\nu\).
\end{proof}

\begin{lemma}[Convexity of the pinching cone]
\label{ch09-lem-pinching-cone-convexity}
\source{chow-knopf-ricci-flow:page-0271, chow-knopf-ricci-flow:page-0272}
\dcref{Lemma 9.5}
\group{positive-section-dominates}

At each fixed time \(t\), the subset \(\mathcal K\) of curvature operators defined by the inequalities
\[
\tr\mathbb P\ge -3(1+t),\qquad
\tr\mathbb P\ge |\nu(\mathbb P)|\bigl(\log|\nu(\mathbb P)|+\log(1+t)-3\bigr)
\]
whenever \(\nu(\mathbb P)\le -1/(1+t)\), is invariant under parallel translation and convex in each fiber.
\end{lemma}

\begin{lemma}[Differential inequality for \(\Omega\)]
\label{ch09-lem-omega-derivative}
\source{chow-knopf-ricci-flow:page-0272, chow-knopf-ricci-flow:page-0273}
\dcref{Lemma 9.6}
\group{positive-section-dominates}

Given a curvature operator \(\mathbb P\) with \(\nu(\mathbb P)<0\), define
\[
\Omega(\mathbb P)=\frac{\tr\mathbb P}{-\nu(\mathbb P)}-\log\!\bigl(-\nu(\mathbb P)\bigr).
\]
Then along the curvature operator of a Ricci flow solution, wherever \(\Omega\) is defined one has
\[
\frac{d}{dt}\Omega\ge -\nu.
\]
\end{lemma}

\begin{corollary}[Limits of Type I and Type IIa singularities]
\label{ch09-cor-nonnegative-curvature-limits}
\source{chow-knopf-ricci-flow:page-0274}
\dcref{Corollary 9.7, Corollary 9.8}
\group{positive-section-dominates}

In dimension \(3\), every Type I limit of a Type I singularity and every Type II limit of a Type IIa singularity has nonnegative sectional curvature for as long as it exists. Moreover, any complete ancient \(3\)-dimensional solution with bounded sectional curvature has nonnegative sectional curvature.
\end{corollary}

\section*{4. Necklike points in finite-time singularities}

\begin{definition}[Type I essential and necklike points]
\label{ch09-def-type-i-essential-necklike}
\source{chow-knopf-ricci-flow:page-0275}
\dcref{Section 9.4}
\group{necklike-points-in-finite-time-singularities}

A point \((x,t)\) is a Type I \(c\)-essential point if \(|\Rm(x,t)|\ge c/(T-t)\), and it is \(\delta\)-necklike if there exists a unit \(2\)-form \(\theta\) at \((x,t)\) such that \(|\Rm-R(\theta\otimes\theta)|\le \delta|\Rm|\).
\end{definition}

\begin{theorem}[Necklike points near Type I and Type IIa singularities]
\label{ch09-thm-necklike-points-finite-time}
\source{chow-knopf-ricci-flow:page-0275, chow-knopf-ricci-flow:page-0276, chow-knopf-ricci-flow:page-0277, chow-knopf-ricci-flow:page-0278, chow-knopf-ricci-flow:page-0279, chow-knopf-ricci-flow:page-0280, chow-knopf-ricci-flow:page-0281, chow-knopf-ricci-flow:page-0282, chow-knopf-ricci-flow:page-0283, chow-knopf-ricci-flow:page-0284}
\dcref{Theorem 9.9}
\group{necklike-points-in-finite-time-singularities}

Let \((\mathcal M^3,g(t))\) be a closed solution of the unnormalized Ricci flow on a maximal time interval \(0<t<T<\infty\). If the normalized flow does not converge to a metric of constant positive sectional curvature, then there exists \(c>0\) such that for all \(\tau\in[0,T)\) and \(\delta>0\), there are \(x\in\mathcal M^3\) and \(t\in[\tau,T)\) such that \((x,t)\) is a Type I \(c\)-essential point and a \(\delta\)-necklike point.
\end{theorem}

\begin{proof}
\source{chow-knopf-ricci-flow:page-0276, chow-knopf-ricci-flow:page-0277, chow-knopf-ricci-flow:page-0278, chow-knopf-ricci-flow:page-0279, chow-knopf-ricci-flow:page-0280, chow-knopf-ricci-flow:page-0281, chow-knopf-ricci-flow:page-0282, chow-knopf-ricci-flow:page-0283, chow-knopf-ricci-flow:page-0284}
\sketch The source modifies the pinching argument from Chapter 6 by replacing \(R\) with \(R+\rho\) and inserting a time factor \((T-t)^\varepsilon\) into the pinching quantity \(F\). Lemma 9.10 gives a comparison of \(|\Rm|\) with \(R+\rho\), Lemma 9.11 provides the evolution inequality for \(F\), Lemmas 9.12 and 9.13 show that \(F\to 0\) when there are no essential necklike points, and Perelman's injectivity estimate supplies the needed noncollapsing control. The resulting limit of dilations is a complete ancient \(3\)-dimensional solution with bounded nonnegative sectional curvature; the absence of necklike points forces positive curvature, and then the pinching limit shows that the limit is a round spherical space form. This contradicts the hypothesis unless the asserted necklike points exist.
\end{proof}

\begin{lemma}[Curvature versus scalar curvature]
\label{ch09-lem-curvature-vs-scalar}
\source{chow-knopf-ricci-flow:page-0276, chow-knopf-ricci-flow:page-0277}
\dcref{Lemma 9.10}
\group{necklike-points-in-finite-time-singularities}

There exists a constant \(C=C(g(0))\) such that \(|\Rm|\le C(R+\rho)\).
\end{lemma}

\begin{lemma}[Evolution of the modified pinching quantity]
\label{ch09-lem-modified-pinching-evolution}
\source{chow-knopf-ricci-flow:page-0277, chow-knopf-ricci-flow:page-0278}
\dcref{Lemma 9.11}
\group{necklike-points-in-finite-time-singularities}

The quantity
\[
F=(T-t)^\varepsilon \frac{|\Rm|^2}{(R+\rho)^{2-\varepsilon}}
\]
satisfies a differential inequality of the form
\[
\frac{\partial}{\partial t}F\le \Delta F+\frac{2(1-\varepsilon)}{R+\rho}\langle \nabla F,\nabla R\rangle+H,
\]
with \(H\) given by the source and decomposed into bad terms \(B_1,B_2\) and good terms \(G_1,G_2\).
\end{lemma}

\begin{lemma}[Pinching improvement away from necks]
\label{ch09-lem-pinching-improves-away-from-necks}
\source{chow-knopf-ricci-flow:page-0278, chow-knopf-ricci-flow:page-0279, chow-knopf-ricci-flow:page-0280}
\dcref{Lemma 9.12, Lemma 9.13}
\group{necklike-points-in-finite-time-singularities}

If \(|\Rm-R(\theta\otimes\theta)|^2\ge \delta |\Rm|^2\) for some \(\delta\in(0,1)\), then \(P\ge \frac{\delta}{96(3-\delta)}|\Rm|^2|\Rm|^{\circ}{}^2\). Consequently, if for every \(c>0\) there are \(\tau\) and \(\delta\) such that after time \(\tau\) every point satisfies either \(|\Rm|<c/(T-t)\) or the above neck-avoidance condition, then \(F\to 0\) uniformly as \(t\to T\).
\end{lemma}

\begin{remark}[Alternate boundedness argument for \(F\)]
\label{ch09-rem-alternate-boundedness-of-F}
\source{chow-knopf-ricci-flow:page-0282, chow-knopf-ricci-flow:page-0283, chow-knopf-ricci-flow:page-0284}
\dcref{Remark 9.14}
\group{necklike-points-in-finite-time-singularities}

The source records an alternate proof that \(F\) remains bounded, using the enhanced maximum principle for systems and a directly controlled auxiliary quantity built from \(\lambda-\nu\).
\end{remark}

\section*{5. Necklike points in ancient solutions}

\begin{lemma}[Ancient solutions have nonnegative scalar curvature]
\label{ch09-lem-ancient-nonnegative-scalar}
\source{chow-knopf-ricci-flow:page-0284}
\dcref{Lemma 9.15}
\group{necklike-points-in-ancient-solutions}

Let \((\mathcal M^n,g(t))\) be a complete ancient solution of the Ricci flow. Assume that \(R_{\min}(t)\) is finite for all \(t\le 0\) and that there is a continuous function \(K(t)\) such that \(|\sect[g(t)]|\le K(t)\). Then \(g(t)\) has nonnegative scalar curvature for as long as it exists.
\end{lemma}

\begin{proof}
\source{chow-knopf-ricci-flow:page-0284}
\sketch The scalar curvature satisfies \(\partial_t R=\Delta R+2|\Ric|^2\ge \Delta R+\frac{2}{n}R^2\). If \(R\) ever becomes nonnegative, it stays so. If \(R_{\min}(t_1)<0\), then the maximum principle gives a lower differential inequality for \(R_{\min}\) that improves as \(t_0\to-\infty\), forcing \(R_{\min}(t)\ge 0\) for all times.
\end{proof}

\begin{remark}[Dimension three]
\label{ch09-rem-ancient-scalar-three-dimensional}
\source{chow-knopf-ricci-flow:page-0284}
\dcref{Remark 9.16}
\group{necklike-points-in-ancient-solutions}

In dimension \(3\), the nonnegative scalar curvature conclusion can be strengthened as in Corollary 9.8.
\end{remark}

\begin{lemma}[Ancient lower blow-up rate]
\label{ch09-lem-ancient-lower-blowup-rate}
\source{chow-knopf-ricci-flow:page-0284, chow-knopf-ricci-flow:page-0285}
\dcref{Lemma 9.17}
\group{necklike-points-in-ancient-solutions}

If \((\mathcal M^n,g(t))\) is an ancient solution with nonnegative Ricci curvature, then there exists \(c_0>0\), depending only on \(n\), such that
\[
\liminf_{t\to -\infty}|t|\sup_{x\in\mathcal M^n}|\Rm(x,t)|\ge c_0.
\]
\end{lemma}

\begin{proof}
\source{chow-knopf-ricci-flow:page-0285}
\sketch The scalar curvature obeys a differential inequality that bounds \((R_{\max})^{-1}\) from below by an affine function of time. Sending \(t\to-\infty\) yields the stated \(\liminf\) lower bound, and the comparison between \(R\) and \(|\Rm|\) finishes the argument.
\end{proof}

\begin{theorem}[Ancient Type I necklike points]
\label{ch09-thm-necklike-points-ancient}
\source{chow-knopf-ricci-flow:page-0285, chow-knopf-ricci-flow:page-0286, chow-knopf-ricci-flow:page-0287}
\dcref{Theorem 9.19}
\group{necklike-points-in-ancient-solutions}

Let \((\mathcal M^3,g(t))\) be a complete ancient solution of the Ricci flow with positive sectional curvature. Suppose that
\[
\sup_{\mathcal M^3\times(-\infty,0]}|t|^\gamma R(x,t)<\infty
\]
for some \(\gamma>0\). Then either \((\mathcal M^3,g(t))\) is isometric to a spherical space form, or else there exists \(c>0\) such that for all \(\tau\in(-\infty,0]\) and \(\delta>0\), there are \(x\in\mathcal M^3\) and \(t\in(-\infty,\tau)\) such that \((x,t)\) is an ancient Type I \(c\)-essential point and a \(\delta\)-necklike point.
\end{theorem}

\begin{proof}
\source{chow-knopf-ricci-flow:page-0285, chow-knopf-ricci-flow:page-0286, chow-knopf-ricci-flow:page-0287}
\sketch The argument is the ancient analogue of Theorem 9.9. One defines
\[
G=|t|^{\gamma\varepsilon/2}\frac{|\Rm|^2}{R^{2-\varepsilon}},
\]
derives a heat-type differential inequality for \(G\), and uses the absence of ancient essential necklike points to force \(G\) to be a subsolution with \(G\to 0\) as \(t\to-\infty\). This implies the solution is locally round, hence globally a spherical space form. Otherwise the asserted necklike points exist.
\end{proof}

\section*{6. Type I ancient solutions on surfaces}

\begin{proposition}[LYH trace estimate]
\label{ch09-prop-lyh-trace-estimate}
\source{chow-knopf-ricci-flow:page-0287}
\dcref{Proposition 9.20}
\group{type-i-ancient-surfaces}

If \((\mathcal M^n,g(t))\) is a solution of the unnormalized Ricci flow on a compact manifold with initially positive curvature operator, then for any vector field \(X\) and all times \(t>0\) for which the solution exists,
\[
0 \le \frac{\partial}{\partial t}R + \frac{R}{t} + 2\langle \nabla R, X\rangle + 2\Rc(X,X).
\]
\end{proposition}

\begin{corollary}[Monotonicity of \(tR\) and \(R\)]
\label{ch09-cor-tR-monotone}
\source{chow-knopf-ricci-flow:page-0287, chow-knopf-ricci-flow:page-0288}
\dcref{Corollary 9.21}
\group{type-i-ancient-surfaces}

If \((\mathcal M^n,g(t))\) is a solution of the unnormalized Ricci flow on a compact manifold with initially positive curvature operator, then \(tR\) is pointwise nondecreasing for all \(t\ge 0\) for which the solution exists. If the solution is also ancient, then \(R\) itself is pointwise nondecreasing.
\end{corollary}

\begin{lemma}[Compactness of complete ancient Type I surfaces]
\label{ch09-lem-ancient-type-i-surface-compact}
\source{chow-knopf-ricci-flow:page-0288}
\dcref{Lemma 9.22}
\group{type-i-ancient-surfaces}

If \((\mathcal N^2,h(t))\) is a complete ancient Type I solution of the Ricci flow with strictly positive scalar curvature, then \(\mathcal N^2\) is compact.
\end{lemma}

\begin{proposition}[Classification of complete ancient Type I surface solutions]
\label{ch09-prop-type-i-surface-classification}
\source{chow-knopf-ricci-flow:page-0288, chow-knopf-ricci-flow:page-0289, chow-knopf-ricci-flow:page-0290}
\dcref{Proposition 9.23}
\group{type-i-ancient-surfaces}

A complete ancient Type I solution \((\mathcal N^2,h(t))\) of the Ricci flow on a surface is a quotient of either a shrinking round \(S^2\) or else a flat \(\mathbb R^2\).
\end{proposition}

\begin{proof}
\source{chow-knopf-ricci-flow:page-0288, chow-knopf-ricci-flow:page-0289, chow-knopf-ricci-flow:page-0290}
\sketch The proof splits according to whether the ancient Type I surface solution has zero or positive scalar curvature. In the positive case, Lemma 9.22 gives compactness, so the surface is \(S^2\) or \(\mathbb{RP}^2\), and after passing to a cover one gets \(S^2\). The entropy is monotone and bounded below, so backward dilations along scalar-curvature maximizers converge to a compact ancient limit of constant entropy. By the classification result cited in the source, that limit is a shrinking round sphere. Equality of entropy then forces the original solution itself to be a shrinking round sphere. In the flat case, the strong maximum principle gives a flat quotient of \(\mathbb R^2\).
\end{proof}

\begin{proposition}[Type II ancient surface limits]
\label{ch09-prop-type-ii-ancient-surface-limits}
\source{chow-knopf-ricci-flow:page-0290}
\dcref{Proposition 9.24}
\group{type-i-ancient-surfaces}

Let \((\mathcal M^2,g(t))\) be a complete Type II ancient solution of the Ricci flow defined on an interval \((-\infty,\omega)\) with \(\omega>0\). Assume there exists a function \(K(t)\) such that \(|R|\le K(t)\). Then either \(g(t)\) is flat or else there exists a backward limit that is the self-similar cigar solution.
\end{proposition}

\begin{proof}
\source{chow-knopf-ricci-flow:page-0290}
\sketch Nonnegative scalar curvature follows from Lemma 9.15, and the strong maximum principle splits off the flat case. Otherwise the scalar curvature is positive and pointwise nondecreasing, so one gets a uniform curvature bound on past time intervals. The injectivity-radius estimate then allows a backward blow-up limit. The Type II hypothesis forces the limit to be eternal, and the cited classification of eternal positively curved surface limits identifies it with the cigar soliton.
\end{proof}

\begin{corollary}[Classification of complete ancient surface solutions]
\label{ch09-cor-ancient-surface-classification}
\source{chow-knopf-ricci-flow:page-0290}
\dcref{Corollary 9.25}
\group{type-i-ancient-surfaces}

If \((\mathcal M^2,g(t))\) is a complete ancient solution defined on \((-\infty,\omega)\), where \(\omega>0\), and its curvature is bounded by some function of time alone, then either the solution is flat, or it is a round shrinking sphere, or there exists a backward limit that is the cigar.
\end{corollary}

\begin{remark}[Chapter references]
\label{ch09-rem-notes-and-commentary}
\dcref{Notes 9.26}
\group{notes-and-commentary}
\source{chow-knopf-ricci-flow:page-0291}
The main references for this chapter are the cited papers on Type I singularities and ancient solutions. The source notes that Theorem 9.9 is essentially Theorem 24.7 of the cited reference, that Theorem 9.19 is analogous to Corollary B3.5 in another cited work, and that Proposition 9.23 is Theorem 26.1 of the cited reference.
\end{remark}
